%% kapa: exam=1 round=36 year=2025 period=2 form=A subject=accounting subject_ko=회계학 seq=6 number=46
%% source: Q-Net 감정평가사 기출문제 (cst00309) artlSeq=5231525 file=raw/1008/5231525/2025년 제36회 감정평가사 1차시험 2교시 문제지 원본.pdf
%% conversion: kordoc 4.24.51 -> tools/gen1.py (auto); review=proofread-sonnet; flags=none
\begin{question}{46}
(주)감평은 20×1년 초 구축물로 분류되는 폐기물처리시설(내용연수 10년, 잔존가치 ￦0, 정액법 상각, 원가모형 적용)을 동 일자에 수령한 정부보조금(상환의무 없음) ￦300,000을 포함하여 총 ￦1,000,000에 취득하였다. 동 시설은 내용연수 종료시점에 원상복구의무가 있으며, 복구시점의 복구비용은 ￦200,000이 소요될 것으로 예상된다. 이는 충당부채의 인식요건을 충족하며, 복구충당부채에 대한 할인율은 연 8 \%이다. 정부보조금과 관련하여 자산차감법으로 인식할 경우, (주)감평이 20×1년에 동 폐기물처리시설과 관련하여 인식할 감가상각비는? (단, 8 \%, 10기간 단일금액 ￦1의 현가계수는 0.4632이고, 화폐금액은 소수점 첫째자리에서 반올림하며, 단수차이로 인한 오차가 있으면 가장 근사치를 선택한다.) \par
\begin{choices}
  \item ￦70,000
  \item ￦79,264
  \item ￦82,562
  \item ￦100,000
  \item ￦109,264
\end{choices}
\end{question}
\answer{2}
